\section*{Chapter \ref{ch:b1:arith} --- Integer Arithmetic}

\begin{solution}{exo:b1:arith:1}
$1001 = 1 \times 777 + 224$; $777 = 3 \times 224 + 105$; $224 = 2
\times 105 + 14$; $105 = 7 \times 14 + 7$; $14 = 2 \times 7 + 0$. So
$\gcd(1001, 777) = 7$. Backwards:
\[
7 = 105 - 7 \times 14
= 105 - 7(224 - 2\times 105) = 15 \times 105 - 7 \times 224
\]
\[
= 15(777 - 3\times 224) - 7\times 224 = 15 \times 777 - 52 \times 224
= 15 \times 777 - 52(1001 - 777) = 67 \times 777 - 52 \times 1001 .
\]
Check: $67 \times 777 = 52\,059$ and $52 \times 1001 = 52\,052$;
difference $7$. Bézout pair: $(u, v) = (-52, 67)$ for $1001u + 777v =
7$.
\end{solution}

\begin{solution}{exo:b1:arith:2}
Since $10 \equiv 1 \pmod 9$: $10^k \equiv 1$, so $\sum_k d_k 10^k
\equiv \sum_k d_k \pmod 9$. Since $10 \equiv -1 \pmod{11}$: $10^k
\equiv (-1)^k$, so the integer is congruent to the alternating sum
$\sum_k (-1)^k d_k$ mod $11$ (starting from the \emph{units} digit
with sign $+$).

$123\,456\,789$: digit sum $45 \equiv 0 \pmod 9$. Alternating sum from
the units: $9 - 8 + 7 - 6 + 5 - 4 + 3 - 2 + 1 = 5$, so the number is
$\equiv 5 \pmod{11}$.
\end{solution}

\begin{solution}{exo:b1:arith:3}
Euclid: $237 = 2 \times 91 + 55$; $91 = 1 \times 55 + 36$; $55 = 1
\times 36 + 19$; $36 = 1 \times 19 + 17$; $19 = 1 \times 17 + 2$; $17
= 8 \times 2 + 1$. Backwards:
\[
1 = 17 - 8\times 2 = 17 - 8(19 - 17) = 9\times 17 - 8\times 19
= 9(36 - 19) - 8\times 19 = 9\times 36 - 17\times 19
\]
\[
= 9\times 36 - 17(55 - 36) = 26\times 36 - 17\times 55
= 26(91 - 55) - 17\times 55 = 26\times 91 - 43\times 55
\]
\[
= 26\times 91 - 43(237 - 2\times 91) = 112 \times 91 - 43 \times 237.
\]
So $91 \times 112 \equiv 1 \pmod{237}$: the solutions are $x \equiv
112 \pmod{237}$. (Check: $91 \times 112 = 10\,192 = 43 \times 237 +
1$.)
\end{solution}

\begin{solution}{exo:b1:arith:4}
$\gcd(17, 39) = 1$: Euclid gives $39 = 2\times 17 + 5$, $17 = 3\times
5 + 2$, $5 = 2\times 2 + 1$, and backwards
\[
1 = 5 - 2\times 2 = 5 - 2(17 - 3\times 5) = 7\times 5 - 2\times 17
= 7(39 - 2\times 17) - 2\times 17 = 7\times 39 - 16\times 17 .
\]
Particular solution $(x_0, y_0) = (-16, 7)$. General solution of the
homogeneous equation $17x + 39y = 0$: $x = 39k$, $y = -17k$ (since
$17 \mid 39y$ and $\gcd(17,39) = 1$ force $17 \mid y$ ---
Gauss's lemma). Hence
\[
(x, y) = (-16 + 39k,\; 7 - 17k), \qquad k \in \Z .
\]
For the right-hand side $5$, multiply the particular solution by $5$:
$(x, y) = (-80 + 39k,\; 35 - 17k)$, $k \in \Z$.
\end{solution}

\begin{solution}{exo:b1:arith:5}
For every prime $p$, with $\alpha = v_p(a)$ and $\beta = v_p(b)$:
\[
v_p\bigl(\gcd(a,b)\bigr) + v_p\bigl(\operatorname{lcm}(a,b)\bigr)
= \min(\alpha, \beta) + \max(\alpha, \beta)
= \alpha + \beta = v_p(ab) .
\]
Two positive integers with the same valuation at every prime are equal
(\cref{prop:b1:arith:valuation}), so $\gcd(a,b)\operatorname{lcm}(a,b)
= ab$.
\end{solution}

\begin{solution}{exo:b1:arith:6}
Valuations: $\gcd(a, b) = 2^{\min(10,6)} 3^{\min(4,7)} 5^{\min(2,0)}
7^{\min(0,1)} = 2^6\, 3^4 = 5184$;
$\operatorname{lcm}(a,b) = 2^{10}\, 3^7\, 5^2\, 7$.

Divisor count: a positive divisor of $n = \prod p_i^{\alpha_i}$ is
exactly a choice $\prod p_i^{\beta_i}$ with $0 \leq \beta_i \leq
\alpha_i$ (\cref{prop:b1:arith:valuation}); the choices are
independent, so there are $\prod_i (\alpha_i + 1)$ divisors. For $a$:
$(10+1)(4+1)(2+1) = 165$.
\end{solution}

\begin{solution}{exo:b1:arith:7}
Suppose $\sqrt p = \frac rq$ with $r, q \in \N^*$, i.e.\ $p q^2 =
r^2$. Apply $v_p$: $v_p(pq^2) = 1 + 2v_p(q)$ is odd, while $v_p(r^2)
= 2 v_p(r)$ is even. An integer cannot have an odd and an even
$p$-valuation at once: contradiction. So $\sqrt p \notin \Q$.
\end{solution}

\begin{solution}{exo:b1:arith:8}
\emph{General fact.} With $\gcd(m,n) = 1$, Bézout gives $mu + nv = 1$.
Set $x_0 = b\,mu + a\,nv$. Then $x_0 \equiv a\,nv \equiv a(1 - mu)
\equiv a \pmod m$ and similarly $x_0 \equiv b \pmod n$: existence. If
$x$ and $x'$ are two solutions, $m$ and $n$ divide $x - x'$, so
$mn \mid x - x'$ (\cref{thm:b1:arith:gauss}~(2)): uniqueness mod
$mn$.

\emph{Numerically:} $m = 7$, $n = 11$: $7 \times (-3) + 11 \times 2 =
1$. So $x_0 = 5 \times 7 \times (-3) + 2 \times 11 \times 2 = -105 +
44 = -61 \equiv 16 \pmod{77}$. Check: $16 = 2\times 7 + 2 \equiv 2
\pmod 7$; $16 = 11 + 5 \equiv 5 \pmod{11}$. Solutions: $x \equiv 16
\pmod{77}$.
\end{solution}

\begin{solution}{exo:b1:arith:9}
Mod $7$: Fermat gives $3^6 \equiv 1$, and $1000 = 6 \times 166 + 4$,
so $3^{1000} \equiv 3^4 = 81 \equiv 4 \pmod 7$.

Last two digits of $7^{100}$: work mod $4$ and mod $25$. Mod $4$: $7
\equiv -1$, so $7^{100} \equiv 1$. Mod $25$: $7^2 = 49 \equiv -1$, so
$7^4 \equiv 1$ and $7^{100} = (7^4)^{25} \equiv 1$. By the Chinese
remainder theorem (\cref{exo:b1:arith:8}), $7^{100} \equiv 1
\pmod{100}$: the last two digits are $01$.
\end{solution}

\begin{solution}{exo:b1:arith:10}
Write $m = nq + r$, $0 \leq r < n$. Then
\[
2^m - 1 = 2^r\bigl(2^{nq} - 1\bigr) + 2^r - 1,
\]
and $2^n - 1$ divides $2^{nq} - 1 = (2^n - 1)(2^{n(q-1)} + \dots +
1)$. So mod $2^n - 1$, $\;2^m - 1 \equiv 2^r - 1$, and since $0 \leq
2^r - 1 < 2^n - 1$, this \emph{is} the Euclidean remainder.

Therefore the Euclidean algorithm on the pair $(2^m - 1, 2^n - 1)$
mirrors, exponent by exponent, the algorithm on $(m, n)$: each
division step replaces $(m, n)$ by $(n, r)$ upstairs and $(2^m - 1,
2^n - 1)$ by $(2^n - 1, 2^r - 1)$ downstairs. The algorithm upstairs
terminates at $\gcd(m,n)$, so downstairs it terminates at
$2^{\gcd(m,n)} - 1$.
\end{solution}

\begin{solution}{exo:b1:arith:11}
First solve $x^2 \equiv 1 \pmod p$: $p \mid (x-1)(x+1)$, so by
Euclid's lemma $x \equiv 1$ or $x \equiv -1 \pmod p$.

In the product $(p-1)! = 1 \times 2 \times \dots \times (p-1)$, every
factor $a$ is invertible mod $p$, and its inverse $a^{-1}$ is again
one of the factors (\cref{prop:b1:arith:invmod}). Pair each $a$ with
$a^{-1}$: the pairs multiply to $1$, except that self-paired factors
($a = a^{-1}$, i.e.\ $a^2 \equiv 1$) stand alone --- and these are
exactly $1$ and $p - 1$. Hence
\[
(p-1)! \equiv 1 \times (p - 1) \equiv -1 \pmod p .
\]
(For $p = 2$: $1! = 1 \equiv -1 \pmod 2$; the pairing argument
degenerates but the result holds.)

\emph{Converse.} Let $n \geq 2$ be composite, $n = ab$ with $1 < a
\leq b < n$. If $a < b$, both appear as distinct factors of $(n-1)!$,
so $n \mid (n-1)!$ and $(n-1)! \equiv 0 \not\equiv -1$. If $a = b$
(i.e.\ $n = a^2$): for $a \geq 3$, both $a$ and $2a$ are $< n$, so
$n = a^2 \mid a \times 2a \mid (n-1)!$, same conclusion; for $n = 4$,
$(n-1)! = 6 \equiv 2 \not\equiv -1 \pmod 4$.
\end{solution}

\begin{solution}{exo:b1:arith:12}
\begin{enumerate}
  \item Induction. For $n = 1$: $F_0 = 3 = F_1 - 2 = 5 - 2$.
        Assuming $F_0\cdots F_{n-1} = F_n - 2$:
        \[
        F_0 \cdots F_n = (F_n - 2)F_n
        = \bigl(2^{2^n} - 1\bigr)\bigl(2^{2^n} + 1\bigr)
        = 2^{2^{n+1}} - 1 = F_{n+1} - 2 .
        \]
  \item Let $m < n$ and $d = \gcd(F_m, F_n)$. By (1), $F_m$
        divides $F_n - 2$, so $d$ divides both $F_n$ and $F_n -
        2$, hence divides $2$. But every Fermat number is odd, so
        $d = 1$.
  \item Each $F_n \geq 3$ has a prime divisor $p_n$
        (\cref{thm:b1:arith:euclidprimes}'s first step). If $m
        \neq n$, then $p_m \neq p_n$, since a common prime would
        divide $\gcd(F_m, F_n) = 1$. The map $n \mapsto p_n$ is
        therefore injective from $\N$ into the primes: there are
        infinitely many primes.
\end{enumerate}
\end{solution}

\begin{solution}{pb:b1:arith:1}
\textbf{1.} $10! = 3\,628\,800$: two trailing zeros. Valuations
factor by factor: powers of $2$ come from $2, 4 = 2^2, 6, 8 = 2^3,
10$, totaling $v_2(10!) = 1 + 2 + 1 + 3 + 1 = 8$; powers of $5$
from $5$ and $10$: $v_5(10!) = 2$. Trailing zeros $=
\min(v_2, v_5) = 2$, consistent.

\textbf{2.} Write the Euclidean division $\lfloor x\rfloor = nq +
r$, $0 \leq r \leq n - 1$. Then $x = nq + r + \{x\}$ with $0 \leq
r + \{x\} < n$, so $\lfloor x/n \rfloor = q = \bigl\lfloor \lfloor
x \rfloor / n \bigr\rfloor$.

\textbf{3.} Let $\alpha = v_p(a) \leq \beta = v_p(b)$ (swap if
needed) and write $a = p^\alpha a'$, $b = p^\beta b'$ with $p
\nmid a', b'$. Then $a + b = p^\alpha\bigl(a' + p^{\beta -
\alpha}b'\bigr)$, so $v_p(a + b) \geq \alpha = \min$. If $\alpha <
\beta$, the parenthesis is $a' + p^{\beta-\alpha}b' \equiv a'
\not\equiv 0 \pmod p$: the valuation is exactly $\alpha$.

\textbf{4.} The multiples of $m$ in $\intint1n$ are $m, 2m, \dots,
qm$ where $q$ is the largest integer with $qm \leq n$, i.e.\ $q =
\lfloor n/m \rfloor$.

\textbf{5.} By unique factorization, $v_p(n!) = \sum_{j=1}^{n}
v_p(j)$. Count differently: each $j$ contributes $v_p(j) =
\#\{k \geq 1 : p^k \mid j\}$, so
\[
v_p(n!) = \sum_{j=1}^n \#\{k : p^k \mid j\}
= \sum_{k\geq1} \#\{j \leq n : p^k \mid j\}
= \sum_{k\geq1} \Bigl\lfloor \frac n{p^k} \Bigr\rfloor
\]
by question 4 --- Legendre's formula. The sum is finite: terms
with $p^k > n$ vanish.

\textbf{6.} $v_5(1000!) = 200 + 40 + 8 + 1 = 249$ (divisions by
$5, 25, 125, 625$); $v_2(1000!) = 500 + 250 + 125 + 62 + 31 + 15 +
7 + 3 + 1 = 994$. Trailing zeros of $1000!$: each zero consumes
one $2$ and one $5$, so there are $\min(994, 249) = 249$ of them.

\textbf{7.} With $n = \sum_i a_ip^i$, question 2 gives $\lfloor
n/p^k \rfloor = \sum_{i \geq k} a_ip^{i-k}$ (truncate the base-$p$
expansion). Summing over $k \geq 1$ and exchanging the two finite
sums:
\[
v_p(n!) = \sum_{i\geq1} a_i \sum_{k=1}^{i} p^{i-k}
= \sum_{i\geq0} a_i\,\frac{p^i - 1}{p - 1}
= \frac{n - s_p(n)}{p - 1} .
\]

\textbf{8.} For $p = 2$: $v_2(n!) = n - s_2(n)$. Since $n \geq 1$
has $s_2(n) \geq 1$, always $v_2(n!) \leq n - 1 < n$: $2^n \nmid
n!$. And $v_2(n!) = n - 1$ iff $s_2(n) = 1$ iff $n$ is a power of
$2$.

\textbf{9.} $n$ has $\lfloor \log_p n \rfloor + 1$ base-$p$
digits, each at most $p - 1$, so $1 \leq s_p(n) \leq
(p-1)\bigl(\log_p(n) + 1\bigr)$. Substituting in question 7:
\[
\frac n{p-1} - \log_p(n) - 1 \;\leq\; v_p(n!) \;<\; \frac n{p-1},
\]
and dividing by $n$: $\frac{v_p(n!)}n \to \frac1{p-1}$.

\textbf{10.} $Z(n) - Z(n-1) = v_5(n!/(n-1)!) = v_5(n)$: the count
of trailing zeros jumps by $v_5(n)$ at each multiple of $5$ and is
constant in between. $Z(24) = \lfloor24/5\rfloor = 4$ and $Z(25) =
5 + 1 = 6$: at $n = 25$ the count jumps from $4$ straight to $6$
($v_5(25) = 2$), and since $Z$ is nondecreasing with $Z \leq 4$
before and $Z \geq 6$ after, the value $5$ is never attained: no
factorial ends in exactly five zeros.

\textbf{11.} Write $x = \lfloor x\rfloor + \{x\}$: $\lfloor x +
y\rfloor = \lfloor x\rfloor + \lfloor y\rfloor + \lfloor \{x\} +
\{y\}\rfloor$, and $0 \leq \{x\} + \{y\} < 2$ makes the last floor
$0$ or $1$. Then, by Legendre applied three times,
\[
v_p\binom{m+n}m = v_p\bigl((m{+}n)!\bigr) - v_p(m!) - v_p(n!)
= \sum_{k\geq1}\Bigl(
\Bigl\lfloor\frac{m+n}{p^k}\Bigr\rfloor
- \Bigl\lfloor\frac{m}{p^k}\Bigr\rfloor
- \Bigl\lfloor\frac{n}{p^k}\Bigr\rfloor\Bigr),
\]
a finite sum of $0$s and $1$s (apply the first claim to $x =
m/p^k$, $y = n/p^k$).

\textbf{12.} Fix $k \geq 1$ and write $m = p^km_1 + m_0$, $n =
p^kn_1 + n_0$ with $0 \leq m_0, n_0 < p^k$ (Euclidean division:
$m_0$ is the number formed by the $k$ low digits of $m$). Then
\[
\Bigl\lfloor\frac{m+n}{p^k}\Bigr\rfloor
- \Bigl\lfloor\frac m{p^k}\Bigr\rfloor
- \Bigl\lfloor\frac n{p^k}\Bigr\rfloor
= \Bigl\lfloor\frac{m_0 + n_0}{p^k}\Bigr\rfloor ,
\]
which is $1$ if $m_0 + n_0 \geq p^k$ and $0$ otherwise. But $m_0 +
n_0 \geq p^k$ says precisely that adding the $k$ low digits of $m$
and $n$ overflows into position $k$ --- a carry into position $k$
in the schoolbook addition algorithm. Summing over $k$:
$v_p\binom{m+n}m$ is the number of carries in the base-$p$
addition $m + n$. (Kummer, 1852.)

\textbf{13.} Apply Kummer to $m = j$, $n = p^k - j$, sum $p^k =
(1\underbrace{0\cdots0}_{k})_p$. Let $a = v_p(j)$, so the base-$p$
digits of $j$ at positions $0, \dots, a-1$ are $0$ and the digit
at position $a$ is nonzero. The digits of $p^k - j$ below position
$a$ are $0$ too ($p^k - j = p^a(p^{k-a} - j/p^a)$). At position
$a$, the two nonzero digits must sum to $p$ (result digit $0$):
one carry; at each position $a+1, \dots, k-1$, digits plus the
incoming carry sum to $p$ (result digit $0$ again): the carry
propagates. Total: $k - a$ carries, so $v_p\binom{p^k}j = k -
v_p(j)$. For $k = 1$: $v_p\binom pj = 1$ for $0 < j < p$, the
divisibility used in \cref{thm:b1:arith:fermat}.

\textbf{14.} By the digital form (question 7), using $s_2(2n) =
s_2(n)$ (appending a zero digit):
\[
v_2\binom{2n}n = \bigl(2n - s_2(2n)\bigr) - 2\bigl(n -
s_2(n)\bigr) = 2s_2(n) - s_2(2n) = s_2(n) \geq 1 :
\]
$\binom{2n}n$ is always even, and $v_2 = 1$ (i.e.\ $\binom{2n}n
\equiv 2 \pmod 4$) exactly when $s_2(n) = 1$, i.e.\ when $n$ is a
power of $2$.

\textbf{15.} Vandermonde with $m = n = k = p$: $\binom{2p}p =
\sum_{j=0}^p \binom pj\binom p{p-j} = \sum_{j=0}^p \binom pj^2$.
For $0 < j < p$, $p \mid \binom pj$ (question 13), so $\binom
pj^2 \equiv 0 \pmod p$; the end terms give $1 + 1$:
$\binom{2p}p \equiv 2 \pmod p$.

\textbf{16.} Base $3$: $500 = 486 + 9 + 3 + 2$, digits (low to
high) $(2, 1, 1, 0, 0, 2)$, so $s_3(500) = 6$; and $1000 = 729 +
243 + 27 + 1$, digits $(1, 0, 0, 1, 0, 1, 1)$, so $s_3(1000) = 4$.
\emph{Kummer:} add $500 + 500$ in base $3$: position $0$: $2 + 2 =
4$, digit $1$ carry $1$; position $1$: $1 + 1 + 1 = 3$, digit $0$
carry $1$; position $2$: $1 + 1 + 1 = 3$, digit $0$ carry $1$;
position $3$: $0 + 0 + 1 = 1$, no carry; position $4$: $0$;
position $5$: $2 + 2 = 4$, digit $1$ carry $1$; position $6$:
carry lands: digit $1$. Four carries: $v_3\binom{1000}{500} = 4$.
\emph{Legendre:} $v_3(1000!) = \frac{1000 - 4}2 = 498$ and
$v_3(500!) = \frac{500 - 6}2 = 247$, so $v_3\binom{1000}{500} =
498 - 2\times247 = 4$. The two computations agree --- and the
addition digits $(1, 0, 0, 1, 0, 1, 1)$ reproduce $1000$, as they
must.

\textbf{17.} By Kummer ($p = 2$, $m = k$, $n' = n - k$):
$\binom nk$ is odd iff the addition $k + (n - k)$ in base $2$ has
no carry, iff at every position the digits satisfy $k_i + (n -
k)_i = n_i$; in that case $k_i \leq n_i$ for all $i$. Conversely,
if $k_i \leq n_i$ for all $i$, then the number with digits $n_i -
k_i$ is $n - k$ and the addition is carry-free. Same proof in base
$p$: $p \nmid \binom nk$ iff every base-$p$ digit of $k$ is at
most the corresponding digit of $n$.

\textbf{18.} Counting the $k \in \intint0n$ whose digits obey
$k_i \leq n_i$: each digit of $k$ is chosen independently among
$n_i + 1$ values, giving $\prod_i (n_i + 1)$ choices; in base $2$
this is $2^{\#\{i : n_i = 1\}} = 2^{s_2(n)}$. Row $4 =
(100)_2$: $2^1 = 2$ odd entries --- indeed $1, 4, 6, 4, 1$ has odd
entries only at the ends. Row $5 = (101)_2$: $2^2 = 4$ --- indeed
$1, 5, 10, 10, 5, 1$.

\textbf{19.} All interior entries even $\iff$ the row has exactly
$2$ odd entries (the two ends always are odd) $\iff 2^{s_2(n)} =
2 \iff s_2(n) = 1 \iff n$ is a power of $2$.

\textbf{20.} In question 11's sum, the $k$-th term vanishes as
soon as $p^k > m + n$ (all three floors are then equal, indeed the
first is $0$ when $p^k > m+n$; more simply each term is $0$).
Hence at most $\lfloor \log_p(m+n)\rfloor$ terms are nonzero, each
worth $1$: $a = v_p\binom{m+n}m \leq \log_p(m+n)$, i.e.\ $p^a \leq
m + n$.

\textbf{21.} For every prime $p$, $v_p\bigl(\operatorname{lcm}(1,
\dots, 2n)\bigr) = \lfloor\log_p(2n)\rfloor$ (the largest power of
$p$ not exceeding $2n$ appears among $1, \dots, 2n$). Question 20
with $m = n$ gives $v_p\binom{2n}n \leq \lfloor\log_p(2n)\rfloor$
for every $p$: by \cref{prop:b1:arith:valuation}, $\binom{2n}n
\mid \operatorname{lcm}(1, \dots, 2n)$. For the size: the ratio
$\binom{2n}{k+1}/\binom{2n}k = \frac{2n-k}{k+1} \geq 1$ exactly
for $k < n$, so the central entry is the largest of the $2n + 1$
entries of row $2n$, whence $4^n = \sum_k \binom{2n}k \leq
(2n+1)\binom{2n}n$. Combining:
\[
\operatorname{lcm}(1, \dots, 2n) \geq \binom{2n}n \geq
\frac{4^n}{2n + 1} .
\]
If there were few primes below $2n$, the lcm could not be this
large: exponential growth of the lcm is a quantitative trace of
the abundance of primes.

\textbf{22.} $Z(n) = \sum_k\lfloor n/5^k\rfloor \approx \frac
n4$, so aim near $n = 4 \times 2026 = 8104$: $Z(8104) = 1620 +
324 + 64 + 12 + 2 = 2022$. Step up by multiples of $5$: $Z(8110)
= 2024$, $Z(8115) = 2025$, and
\[
Z(8120) = 1624 + 324 + 64 + 12 + 2 = 2026 .
\]
Since $Z$ is constant between multiples of $5$ and $Z(8119) =
Z(8115) = 2025$, the smallest $n$ with at least $2026$ trailing
zeros is $n = 8120$.

\textbf{23.} Base $7$: $50 = 49 + 1$, digits (low to high) $(1, 0,
1)$. Adding $50 + 50$: position $0$: $1 + 1 = 2 < 7$, no carry;
position $1$: $0 + 0 = 0$; position $2$: $1 + 1 = 2 < 7$, no
carry. Carry-free, so by Kummer $v_7\binom{100}{50} = 0$: $7 \nmid
\binom{100}{50}$. Legendre agrees: $v_7(100!) = \lfloor 100/7
\rfloor + \lfloor 100/49 \rfloor = 14 + 2 = 16$ and $v_7(50!) = 7
+ 1 = 8$, so $v_7\binom{100}{50} = 16 - 2\times8 = 0$.

\textbf{24.} (i) Unique factorization underlies the very
definition of $v_p$ and its additivity, hence Legendre's formula
and every divisibility conclusion
(\cref{prop:b1:arith:valuation}). (ii) Euclidean division produced
the truncation identity of question 2 and the split $m = p^km_1 +
m_0$ that isolates the carry (question 12). (iii) Counting: the
multiples-of-$m$ count (question 4), the digit-choice product
(question 18), and the row-sum bound $4^n \leq
(2n+1)\binom{2n}n$ (question 21) are all
\cref{ch:b1:counting}-style arguments.

\textbf{25.} Legendre converts ``what power of $p$ divides
$n!$'' into base-$p$ digit arithmetic; Kummer compresses the
answer for binomial coefficients into the carries of a single
addition --- divisibility, seemingly a global property of huge
numbers, is read off locally, digit by digit. Question 16 is the
paradigm: four carries, computed by hand, determine the exact
power of $3$ in a number with hundreds of digits. And question 21
shows the same circle of ideas brushing against deep waters: an
exponential lower bound for $\operatorname{lcm}(1, \dots, 2n)$ is
a first, fully elementary step toward the prime number theorem,
whose proof lies far beyond this volume. The entire toolkit ---
division, gcd, valuations --- is replayed for polynomials in
\cref{ch:b1:poly}, where the analogue of a digit expansion is
expansion in powers of $(X - a)$.
\end{solution}
